\section{Results}
\label{sec:results}

\subsection{Overall Performance}

Figure~\ref{fig:cupath-performance} compares CuPath with Baseline using the
formal 14-workload campaign.  Every pair uses the same frozen binary,
benchmark input, architectural parameters, workgroup-completion rule, and
completed workgroup count; only the CuPath mechanisms differ.  CuPath achieves
a 1.465$\times$ geometric-mean speedup and improves 13 of the 14 workloads.
MM, FWS, and KM obtain the largest gains at 6.171$\times$,
4.155$\times$, and 2.511$\times$, respectively.  BT and FWT improve by
1.491$\times$ and 1.578$\times$, while I2C and RELU improve by
1.214$\times$ and 1.117$\times$.  The remaining workloads stay close to
Baseline, with SPMV at 0.999$\times$.  The result shows that CuPath provides
large gains when requests expose eliminable or aggregatable work, while its
work-conserving fallback keeps workloads with few opportunities near
Baseline.

\begin{figure}[t]
  \centering
  \includegraphics[width=\columnwidth]{../../akkalat/results/2026-07-22-baseline-complete-observation/cupath_baseline_complete_speedup.png}
  \caption{Normalized performance of CuPath relative to Baseline.  Baseline
  is normalized to 1.0 for every workload, and GM is the geometric mean across
  all 14 workloads.}
  \label{fig:cupath-performance}
\end{figure}

\subsection{Removing Work Along the Path}

Figure~\ref{fig:cupath-work} explains the speedup in terms of work that no
longer reaches a downstream component.  Weighted across the fourteen
workloads, Complete removes 22.6\% of L2 tag lookups, 4.6\% of unique remote
reads, 48.2\% of wire packets, and 30.8\% of repeated wafer traversals.  The
different columns peak on different applications: L2 filtering is strongest
for SPMV and MT, packet aggregation for AES, RELU, and KMeans, and requesting-GPM
L2 reuse for MM, FIR, SC, and I2C.  This separation supports the path-oriented
design: no single boundary exposes all four reductions.

CuPath does not obtain these gains by increasing aggregate memory traffic.
Complete reduces physical DRAM reads from 39.25 to 38.71 million (1.38\%) and
writes from 20.83 to 20.68 million (0.71\%).  The sample-weighted local
L2-miss-to-DRAM-issue latency falls from 38.12~ns in Baseline to 20.34~ns over
the ordinary and Filter-negative Complete issue populations.  We report these
populations separately in the machine-readable results; the comparison is not
an assumption that every miss takes the fast path.

\begin{figure}[t]
  \centering
  \includegraphics[width=\columnwidth]{../../akkalat/results/2026-07-17-filter-prefetch-v5-formal14/cupath_work_reduction.png}
  \caption{Percentage of baseline work removed at successive boundaries.
  The bottom row is weighted by the corresponding baseline work count.}
  \label{fig:cupath-work}
\end{figure}

\subsection{M1: Cuckoo-Filter-Guided Local Pairing}

M1 is a supporting transformation rather than the dominant source of
speedup.  It issues 261,615 local candidates over the full suite: 252,850
reach the DRAM-issue point and 8,765 lose a race to an ordinary demand;
176,447 are eventually useful.  Although these issued candidates are
additional logical accesses, exact MSHR merging and demand timing leave total
physical reads 0.35\% lower than Baseline and increase L2 MSHR-full cycles by
only 0.22\%.  FIR is the exception: M1 alone is 0.898$\times$ and increases
physical reads and writes, so we do not claim that local prediction is
universally beneficial.

The coupling experiment isolates why the Filter is still useful for mostly
local workloads.  Filter-only definite-miss elimination provides
1.0145$\times$ geometric mean, while ungated prediction provides only
1.0059$\times$.  Coupling the same predictor to \textsc{Pattern},
\textsc{Resident}, and \textsc{Pending} reaches 1.0188$\times$.  It removes
34.7\% of issued candidates while retaining 96.7\% of useful candidates,
raising useful-per-issued accuracy from 24.0\% to 35.5\%.  Additional DRAM
reads fall from 6,206 to 5,578 and demand-delay exposure events from 145,466
to 33,066.  Predictor-only records candidates but is exactly 1.0000$\times$,
confirming that observation itself is not counted as timing benefit.

Replacing Cuckoo membership with exact metadata changes the same screen only
from 1.0188$\times$ to 1.0192$\times$.  The Cuckoo design observes four false
positives in 273,146 queries (0.0015\%).  Thus the Filter does not move data or
create prediction benefit by itself; it compactly rejects avoidable work and
tracks the request lifecycle with little loss relative to exact metadata.

\subsection{M2: Cuckoo-Filter-Guided Remote Aggregation}

M2 is the largest standalone contributor.  Complete records 4.24 million
real demands merged by exact RDMA line state and 2.52 million predicted lines
piggybacked into existing home-GPM/page packets.  Because a real demand can take
over a candidate before batch insertion, the aggregate counters do not expose
the intersection of these two populations; we therefore do not relabel every
useful prediction as an avoided packet.  The exact merge and wire-packet
counters independently establish the reductions.

\subsection{M3: Cuckoo-Filter-Guided Remote Reuse}

M3 bypasses 23.41 million exact requesting-GPM L2 probes when \textsc{Resident}
and \textsc{Seen} reliably prove absence, 42.6\% of eligible unique-request
probe opportunities.  The remaining probes produce 27.23 million physical
L2 hits and satisfy 28.21 million logical responses through exact waiter
fanout.  M3 installs 7.06 million remote-clean lines in the existing L2;
0.996 million are observed to retire unused, a 14.1\% lower bound because
lines still resident at termination are unclassified.  Invalid ways absorb
5.35 million fills and older remote-clean lines absorb the remaining 1.71
million; no ordinary local-clean line is displaced.  The largest
workload-level sum of per-slice peaks occupies 37.5\% of wafer L2 lines, a
conservative non-simultaneous upper bound rather than an average footprint.
Of the unused retirements, 160,257 carry candidate provenance and remove the
associated \textsc{Pattern}.  A useful patterned hit instead keeps
\textsc{Pattern} eligible; the aggregate counter does not isolate that subset
from all requesting-GPM L2 hits.

\subsection{Sensitivity and Limitations}

Figure~\ref{fig:cupath-sensitivity} varies only the parameters needed to audit
the selected point.  Across four workloads, all tested points span
1.2713$\times$--1.2817$\times$, compared with 1.2757$\times$ for the
32K-slot, 13-bit, 64-entry, one-cycle point.  Reducing the Filter to 16K slots
raises insertion failures from 27 to 156,726, whereas 64K slots, 16
fingerprint bits, or a 256-entry predictor changes geometric mean by less
than 0.5\%.  A zero-cycle lookup reaches 1.2791$\times$, only 0.27\% above the
selected point; the result therefore does not depend on a free Filter lookup.

The long runs expose a limitation hidden by the short screen.  Across the
fourteen Complete runs, correctness-critical \textsc{Resident} and
\textsc{Pending} have zero insertion failures, but capacity pressure causes
9.42 million \textsc{Pattern} and 5.81 million \textsc{Seen} attempts to fail
closed.  These failures suppress optional prediction or reuse and cannot
change demand correctness, but they bound coverage.  The modest
1.187$\times$ Remote-group speedup and M1's FIR regression likewise show that
CuPath is constrained by available aggregation and safe idle opportunities,
not by an unlimited speculative upper bound.

A controlled simulator-runtime screen uses seven representative workloads,
all five configurations, 192 drained workgroups, one host worker, and the
same binary.  Normalized by simulated GPU time, Complete costs
1.190$\times$ Baseline; this is simulator overhead, not hardware performance.

\begin{figure}[t]
  \centering
  \includegraphics[width=\columnwidth]{../../akkalat/results/2026-07-17-filter-prefetch-v5-sensitivity/cupath_sensitivity.png}
  \caption{Four-workload sensitivity.  Orange marks the selected point and
  blue marks alternatives; all use the same frozen binary.}
  \label{fig:cupath-sensitivity}
\end{figure}
